\chapter{Introdução}
\label{chap:introducao}

\section{Contextualização e problema de pesquisa}
\label{sec:problema}

\todo{Apresentar o contexto de propulsão elétrica, o desafio multidisciplinar do projeto de propellers e formular claramente o problema de pesquisa.}

\section{Justificativa}
\label{sec:justificativa}

\todo{Justificar a relevância de uma ferramenta de projeto preliminar que integre critérios aerodinâmicos, estruturais e aeroacústicos com custo computacional reduzido.}

\section{Objetivos}
\label{sec:objetivos}

\subsection{Objetivo geral}

Desenvolver e validar um algoritmo computacional para otimização geométrica de \propellers, considerando critérios estruturais, aerodinâmicos e aeroacústicos por meio de modelos analíticos implementados computacionalmente.

\subsection{Objetivos específicos}

\begin{itemize}[leftmargin=*]
    \item Identificar os principais parâmetros geométricos, operacionais e materiais que influenciam o desempenho de \propellers;
    \item Desenvolver e consolidar o modelo aerodinâmico baseado em Blade Element Momentum Theory (BEMT) para estimativa de empuxo, torque e eficiência;
    \item Definir e implementar uma abordagem estrutural simplificada para avaliação da integridade da pá;
    \item Consolidar o modelo aeroacústico baseado na formulação de Hanson para estimativa de componentes tonais;
    \item Definir uma estratégia de otimização capaz de maximizar o desempenho aerodinâmico, reduzir a emissão acústica e respeitar restrições estruturais;
    \item Validar os módulos e o algoritmo por comparação com dados experimentais e/ou modelos de referência;
    \item Avaliar as limitações da formulação proposta e indicar possibilidades de continuidade.
\end{itemize}

\section{Estrutura do trabalho}
\label{sec:estrutura_trabalho}

\todo{Ao final da redação, resumir em um parágrafo o conteúdo de cada capítulo.}
